% GENERATED FILE. Edit canonical lessons, then run npm run generate:books.
% canonical-source-hash: fnv1a64:11449b33212a8f35
% canonical-lessons: MR-C244-mohari, MR-C244-kothimbir, MR-C244-kadhipatta, MR-C244-sengdane, MR-C244-kaju

\chapter{Mustard seed, Coriander leaves, Curry leaves, Peanuts, Cashews}
\label{ch:mr-mustard-seed-coriander-leaves-curry-leaves-peanuts-cashews}
\hlchaptermodality{244}

\begin{chapteropening}
\textbf{By the end of this chapter:} \emph{I can say mustard seed, coriander leaves, curry leaves, peanuts and cashews.}

Complete the last lesson of chapter 244: I can say mustard seed, coriander leaves, curry leaves, peanuts and cashews.
\end{chapteropening}

\section[\texorpdfstring{\mr{मोहरी} (moharī)}{मोहरी (moharī)}]{\mr{मोहरी} (moharī) — mustard seed}

\label{lesson:MR-C244-mohari}



\begin{quote}
Before the new one: say the Marathi for a griddle, then the Marathi for a wok.
\end{quote}

\subsection*{You'll want to know: \mr{मोहरी}}
\textbf{\mr{मोहरी}} — \emph{moharī} — \textquotedblleft{}mustard seed\textquotedblright{}.

\begin{grammarlens}[title={What you've built}]
One more word for this chapter.
\end{grammarlens}

\subsection*{Your turn}
\begin{itemize}
\raggedright
  \item \emph{Say it:} \emph{moharī}
  \item \emph{Say it:} \emph{moharī}, once more
  \item \emph{Say it:} say \emph{moharī}
  \item \emph{Recall:} say the Marathi for a griddle, then the Marathi for a wok, then say \emph{moharī} again
\end{itemize}

\begin{tcolorbox}[breakable,colback=teal!4,colframe=teal!35!black,title={Before you move on}]
What does \mr{मोहरी} mean? (\textquotedblleft{}Mustard seed\textquotedblright{}.) Say it once more.
\end{tcolorbox}
\section[\texorpdfstring{\mr{कोथिंबीर} (kothimbīr)}{कोथिंबीर (kothimbīr)}]{\mr{कोथिंबीर} (kothimbīr) — coriander leaves}

\label{lesson:MR-C244-kothimbir}



\begin{quote}
Before the new one: say the Marathi for a wok, then the Marathi for mustard seed.
\end{quote}

\subsection*{You'll want to know: \mr{कोथिंबीर}}
\textbf{\mr{कोथिंबीर}} — \emph{kothimbīr} — \textquotedblleft{}coriander leaves\textquotedblright{}.

\begin{grammarlens}[title={What you've built}]
One more word for this chapter.
\end{grammarlens}

\subsection*{Your turn}
\begin{itemize}
\raggedright
  \item \emph{Say it:} \emph{kothimbīr}
  \item \emph{Say it:} \emph{kothimbīr}, once more
  \item \emph{Say it:} say \emph{kothimbīr}
  \item \emph{Recall:} say the Marathi for a wok, then the Marathi for mustard seed, then say \emph{kothimbīr} again
\end{itemize}

\begin{tcolorbox}[breakable,colback=teal!4,colframe=teal!35!black,title={Before you move on}]
What does \mr{कोथिंबीर} mean? (\textquotedblleft{}Coriander leaves\textquotedblright{}.) Say it once more.
\end{tcolorbox}
\section[\texorpdfstring{\mr{कढीपत्ता} (kaḍhīpattā)}{कढीपत्ता (kaḍhīpattā)}]{\mr{कढीपत्ता} (kaḍhīpattā) — curry leaves}

\label{lesson:MR-C244-kadhipatta}



\begin{quote}
Before the new one: say the Marathi for mustard seed, then the Marathi for coriander leaves.
\end{quote}

\subsection*{You'll want to know: \mr{कढीपत्ता}}
\textbf{\mr{कढीपत्ता}} — \emph{kaḍhīpattā} — \textquotedblleft{}curry leaves\textquotedblright{}.

\begin{grammarlens}[title={What you've built}]
One more word for this chapter.
\end{grammarlens}

\subsection*{Your turn}
\begin{itemize}
\raggedright
  \item \emph{Say it:} \emph{kaḍhīpattā}
  \item \emph{Say it:} \emph{kaḍhīpattā}, once more
  \item \emph{Say it:} say \emph{kaḍhīpattā}
  \item \emph{Recall:} say the Marathi for mustard seed, then the Marathi for coriander leaves, then say \emph{kaḍhīpattā} again
\end{itemize}

\begin{tcolorbox}[breakable,colback=teal!4,colframe=teal!35!black,title={Before you move on}]
What does \mr{कढीपत्ता} mean? (\textquotedblleft{}Curry leaves\textquotedblright{}.) Say it once more.
\end{tcolorbox}
\section[\texorpdfstring{\mr{शेंगदाणे} (śeṅgdāṇe)}{शेंगदाणे (śeṅgdāṇe)}]{\mr{शेंगदाणे} (śeṅgdāṇe) — peanuts}

\label{lesson:MR-C244-sengdane}



\begin{quote}
Before the new one: say the Marathi for coriander leaves, then the Marathi for curry leaves.
\end{quote}

\subsection*{You'll want to know: \mr{शेंगदाणे}}
\textbf{\mr{शेंगदाणे}} — \emph{śeṅgdāṇe} — \textquotedblleft{}peanuts\textquotedblright{}.

\begin{grammarlens}[title={What you've built}]
One more word for this chapter.
\end{grammarlens}

\subsection*{Your turn}
\begin{itemize}
\raggedright
  \item \emph{Say it:} \emph{śeṅgdāṇe}
  \item \emph{Say it:} \emph{śeṅgdāṇe}, once more
  \item \emph{Say it:} say \emph{śeṅgdāṇe}
  \item \emph{Recall:} say the Marathi for coriander leaves, then the Marathi for curry leaves, then say \emph{śeṅgdāṇe} again
\end{itemize}

\begin{tcolorbox}[breakable,colback=teal!4,colframe=teal!35!black,title={Before you move on}]
What does \mr{शेंगदाणे} mean? (\textquotedblleft{}Peanuts\textquotedblright{}.) Say it once more.
\end{tcolorbox}
\section[\texorpdfstring{\mr{काजू} (kājū)}{काजू (kājū)}]{\mr{काजू} (kājū) — cashews}

\label{lesson:MR-C244-kaju}



\begin{quote}
Before the new one: say the Marathi for curry leaves, then the Marathi for peanuts.
\end{quote}

\subsection*{You'll want to know: \mr{काजू}}
\textbf{\mr{काजू}} — \emph{kājū} — \textquotedblleft{}cashews\textquotedblright{}.

\begin{grammarlens}[title={What you've built}]
One more word for this chapter.
\end{grammarlens}

\subsection*{Your turn}
\begin{itemize}
\raggedright
  \item \emph{Say it:} \emph{kājū}
  \item \emph{Say it:} \emph{kājū}, once more
  \item \emph{Say it:} say \emph{kājū}
  \item \emph{Recall:} say the Marathi for curry leaves, then the Marathi for peanuts, then say \emph{kājū} again
\end{itemize}

\begin{tcolorbox}[breakable,colback=teal!4,colframe=teal!35!black,title={Before you move on}]
What does \mr{काजू} mean? (\textquotedblleft{}Cashews\textquotedblright{}.) Say it once more.
\end{tcolorbox}
